% Source: https://artofproblemsolving.com/community/c3h1761309_imc_2004_q21
% Stable identity: aops-topic-1761309-q01
\textbf{One solution by the Euclidean algorithm.} $17=5\cdot3+2$ and $3=1\cdot2+1$. Working backwards,
\[
1=3-2=3-(17-5\cdot3)=6\cdot3-1\cdot17,
\]
so $m=-1$, $n=6$ works. (Equivalently: $17m+3n=1$ forces $17m\equiv1\pmod3$, i.e.\ $2m\equiv1$, so $m\equiv2\equiv-1\pmod3$; the choice $m=-1$ gives $3n=18$, $n=6$.)

\textbf{All solutions.} If $17m+3n=1=17(-1)+3\cdot6$, then $17(m+1)=3(6-n)$. Since $\gcd(17,3)=1$, $3\mid m+1$; write $m+1=3k$, and then $6-n=17k$. Conversely every such pair works. So the solutions are
\[
\boxed{(m,n)=(-1+3k,\ 6-17k),\ k\in\mathbb Z}.
\]
